% =====================================================================
%  CH 12-16 : Gouvernance · Menaces et limites · Jalons · Correspondance · Décisions
% =====================================================================
\section{Gouvernance}\label{sec:gouvernance}

\subsection{Journal des décisions — format normalisé}\label{sec:journal}

Le protocole ne se révisant jamais en silence, chaque décision structurante est consignée
dans un journal. La v0.4.0 normalise le format (D24) :

\begin{table}[htbp]
\centering\small
\caption{Format normalisé d'une entrée du journal des décisions.}
\label{tab:formatdecisions}
\begin{tabularx}{\textwidth}{@{}l Y@{}}
\toprule
\textbf{Champ} & \textbf{Exigence} \\
\midrule
ID & \code{DNNN}, séquentiel, jamais réutilisé \\
Date & date de consignation \\
Statut & active · amendée par Dxx · invalidée par Dxx \\
Contexte & exigence, incident ou évaluation qui motive la décision \\
Choix retenu & décision en une phrase, rattachée à la fonction objectif \\
Alternatives rejetées & au moins une, avec motif \\
Amendements & référence explicite à la décision remplacée, le cas échéant \\
Développement & argumentaire détaillé si la décision engage l'architecture \\
\bottomrule
\end{tabularx}
\end{table}

Trois règles : le journal est \textbf{append-only} (une entrée ne s'efface jamais) ; une
révision \textbf{invalide explicitement} ce qu'elle remplace ; aucune entrée n'est
rétrodatée. La conformité du journal — présence des champs, unicité des ID, chaînage des
amendements — est contrôlée mécaniquement par l'audit de gouvernance.

\subsection{Tableau de bord de conformité}\label{sec:tableaudebord}

La v0.4.0 unifie les indicateurs de santé en un \textbf{tableau de bord} (D24), généré par
la commande d'audit à chaque commit et à chaque cristallisation, archivé avec la carte :

\begin{table}[htbp]
\centering\small
\caption{Tableau de bord de conformité : métriques, cibles, sources.}
\label{tab:dashboard}
\begin{tabularx}{\textwidth}{@{}Y l Y@{}}
\toprule
\textbf{Métrique} & \textbf{Cible} & \textbf{Source de mesure} \\
\midrule
Blocs avec spécification I à jour & 100 \% & audit de traçabilité \\
Chaînes de traçabilité complètes I $\leftrightarrow$ C $\leftrightarrow$ B
$\leftrightarrow$ tests & 100 \% & audit de traçabilité \\
Construction & 0 / 0 & build \\
Tests différentiels & 100 \% verts & oracle \\
Empreintes des données embarquées & conformes & hachage vs carte \\
Défauts ouverts par classe T1–T6 & 0 & journal des décisions \\
Optimisations sans mesure (R11) & 0 & journal d'optimisation \\
Budgets tokens (localiser / éditer / ajouter) & $\leq$ 64 / 3 000 / 6 000 &
instrumentation — mesurés M1 : 51 / 719 / 1 088 (D31) \\
Suite de conformité par profil actif & verte & suite du profil \\
Changements hors journal des décisions & 0 & diff vs journal \\
Pré-versions correctement balisées (dev / rc / final) & 100 \% & registre des balises \\
\bottomrule
\end{tabularx}
\end{table}

Le tableau de bord est la contrepartie mesurable de la gouvernance : chaque affirmation du
protocole s'y retrouve sous forme de métrique avec sa cible et sa source — aucune ne
s'appuie sur une déclaration non vérifiable. Les métriques dépendantes de l'outillage
affichent leur statut « cible » jusqu'à leur premier mesurage.

\subsection{Continuité et anti-enfermement}

L'export de continuité (D16) demeure l'unique forme de transfert de l'état de
développement : conteneur chiffré et signé regroupant I, B, C et un manifeste
autodescriptif ; oracle régénéré depuis I à l'arrivée, jamais exporté ; import en cinq
étapes dont l'échec d'une seule est terminal. Sauvegarde 3-2-1 chiffrée et escrow de la
clé maîtresse inchangés.

\paragraph{Le coffre-fort des phrases maîtresses (D75).} Le composant qui centralise
toutes les phrases maîtresses d'un projet — et de tout projet ultérieur conduit sous le
protocole — dans un conteneur unique chiffré TRIADE-CF1 (\code{coffre\_fort.py}) :
phrase de contrôle sous garde durable (D81 — sauvegardée dès production au dépôt et au
racine du coffre Drive, AVANT toute autre action), Argon2id aux
planchers de la carte au repos puis HKDF à contexte dédié, XChaCha20-Poly1305 avec
l'en-tête clair pour AAD, écriture atomique 0600, autotest de falsification exigé à
chaque installation. Le conteneur chiffré s'exporte tel quel vers le coffre Drive du
projet et le coffre personnel du développeur — son vol relève de la menace « vol du
coffre-fort des phrases », étendue en D81. Après une réinitialisation du poste, l'import
du conteneur, la relecture de la phrase depuis \code{PHRASE-CONTROLE-COFFRE.txt}, puis
\code{obtenir --vers-fichier} restaurent la phrase opérationnelle sans cérémonie
de régénération. Discipline évoluant en D81 sur un point, l'arbitrage (risque de PERTE
vécu contre risque d'EXPOSITION théorique) étant consigné : chaque cérémonie de clés
(D27/D30) alimente le coffre dans la même session, la phrase de contrôle est sauvegardée
durablement AVANT TOUTE AUTRE ACTION — \code{PHRASE-CONTROLE-COFFRE.txt} au racine du
dépôt (commit dédié, push GitHub, miroir Drive par le hook) et au racine du coffre
Drive, md5 vérifiés : le fichier est la remise, l'affichage unique cesse d'être le canal
de survie — et la graine Ed25519 reste hors coffre (D30). Contrepartie consignée : la
phrase vit sur le stockage du coffre Drive partagé par lien — la confidentialité au
repos repose sur le contrôle d'accès au dossier (restreindre le partage par lien :
mitigation première) ; une phrase entrée dans l'historique git y demeure à vie, la
rotation reste le remède — format normatif : FORMAT\_CARTE § 11.

\section{Menaces et limites}\label{sec:menaces}

Le principe directeur demeure : séparer l'obfuscation, qui protège la propriété
intellectuelle, du durcissement, qui protège l'exécution. Conformément au principe de
Kerckhoffs, la sécurité ne repose jamais sur le secret de la carte : elle survivrait à sa
divulgation.

\subsection{Modèle de menace (étendu en v0.4.0, D20)}

\begin{table}[htbp]
\centering\small
\caption{Modèle de menace — ajouts v0.4.0 en gras ; ajout v0.5.1 en gras (D75), étendu v0.5.2 (D81) ; ajouts v0.6.0 en gras (D84--D87, D121, D133).}
\label{tab:menaces}
\begin{tabularx}{\textwidth}{@{}Y Y Y@{}}
\toprule
\textbf{Menace} & \textbf{Contre-mesure} & \textbf{Résiduel} \\
\midrule
Rétro-ingénierie de l'IP & carte privée ; binaire sémantiquement muet ; variantes par
build & temps acheté — s'érode avec les LLM de décompilation \\
Empoisonnement de carte & signature ; refus de carte non signée ; verrou croisé (R8, R9)
& détecté — session bloquée \\
Corruption de carte (accidentelle) & AEAD par fragment + signature + verrou & détecté au
chargement \\
Fuite de la carte & chiffrement ; sharding ; escrow & perte limitée à un fragment \\
Vol de l'export de continuité & chiffrement intégral + signature & indéchiffrable sans
la phrase \\
\textbf{Vol du coffre-fort des phrases} & chiffrement intégral (Argon2id, contexte
HKDF dédié, AEAD) ; écriture atomique ; depuis D81 la phrase de contrôle vit au coffre
Drive et au dépôt — mitigation première : restreindre le partage par lien & confidentialité
ramenée au contrôle d'accès du dossier ; une phrase commise en git y demeure à vie — la
rotation reste le remède (D75, étendu D81) \\
\textbf{Expropriation (employeur, tiers)} & barrière juridique BSL 1.1 (D84) ;
barrière cryptographique : clé maîtresse jamais entière (Shamir, axe 2), intention
signée par la clé hors-ligne (sceau I, axe 3), serrure piégée à média séparé (axe 4) &
le code volé est inexploitable et ne peut pas évoluer ; l'axe 1 (TEE) reste une
exigence consignée \\
\textbf{Binaire altéré} & checksums par bloc + signature du manifeste au lancement
(R10 a, b) & détecté — arrêt avant tout effet \\
\textbf{Rollback non autorisé} & monotonie de version signée, plancher anti-rollback
(R10 c) & détecté — refus \\
\textbf{Evil maid sur le conteneur ou le média du piège} & tag du piège, époques
croisées conteneur/piège, AEAD + étiquette, champ \code{verif} indérivable (CF3
§ 5) & détecté — refus nommé ; la falsification prouvée détruit côté machine,
jamais le média \\
\textbf{Roll-back du couple conteneur + piège} & trois témoins d'époque : journal
local, carte papier du sachet (à la main), journal Git poussé (CF3 § 9) & détecté —
« RETOUR EN ARRIÈRE » nommé ; recoupement humain à froid \\
\textbf{Agent défaillant} & R2 ; R5 ; oracle dérivé de I ; budgets & l'erreur ne devient
jamais permanente sans vérification \\
\textbf{Données embarquées modifiées} & empreintes dans la carte ; vérification au
lancement (R10 b) & détecté — arrêt propre \\
\textbf{Environnement non conforme} & contrôles de conformité au profil (R10 d) ;
fail-closed & refus propre avant exécution \\
Fuzzing boîte noire & contrats, oracle, tests par propriétés & indépendant du secret de
la carte \\
Redistribution frauduleuse & empreintes par build — canary trap & fuite traçable \\
Instrumentation dynamique & anti-debug léger par mesure de temps & assumé \\
\bottomrule
\end{tabularx}
\end{table}

Les six menaces en gras sont les additions de la v0.4.0, le vol du coffre-fort (D75)
celle de la v0.5.1 ; toutes sont neutralisées par des mécanismes \textbf{locaux et
vérifiables} — la contre-mesure anti-rollback ne dépend d'aucun service en ligne (le
chargement doit échouer proprement hors ligne), et les vérifications de R10 s'exécutent
avant toute entrée dans le chemin d'exécution.

\subsection{Limites assumées}\label{sec:limites}

La rétro-ingénierie est un objectif d'économie, pas d'impossibilité : le coût de l'analyse
d'un binaire strippé, sans structure, dont chaque exemplaire diffère, reste supérieur à sa
valeur — mais les modèles de langage appliqués à la décompilation progressent, et
l'avantage doit être réévalué à chaque génération. L'anti-debug reste léger et assumé.
La furtivité face aux EDR est hors périmètre : un outil légitime par construction, dont
la transparence vis-à-vis du système est totale. Enfin, la gouvernance du protocole a été
validée sur un cas unique (N = 1) — l'étude de cas RoomEq, dont les limites sont
documentées dans l'édition académique (D18) : la généralité reste à établir par
réplication, le cœur natif par les jalons M0 à M7.

\section{Critères d'acceptation des jalons M0 à M10}\label{sec:jalons}

Onze jalons, chacun avec des critères de sortie mesurables (détail dans
\code{PLAN\_MILESTONES.md}). \textbf{Conditions transverses} (v0.4.0) : aucun jalon n'est
franchi sans oracle vert, suite de conformité verte pour les éléments couverts, tableau de
bord conforme et carte signée.

\begin{table}[htbp]
\centering\small
\caption{Critères d'acceptation des jalons.}
\label{tab:jalons}
\begin{tabularx}{\textwidth}{@{}l Y >{\raggedright\arraybackslash}p{2.6cm}@{}}
\toprule
\textbf{Jalon} & \textbf{Critères de sortie mesurables} & \textbf{Statut} \\
\midrule
M0 — chaîne artefact & exécution sur Windows 11 ; round-trip 100 \% ; carte v0 signée ;
exécutable $<$ 4 Kio ; R1–R3, R8–R9 dès le premier octet & \textbf{franchi} 26/10/2026 (D26) \\
M1 — blocs réutilisables & oracle 100 \% sur la mini-libc ; première cristallisation ;
adresses stables démontrées ; première mesure effective des budgets tokens & \textbf{franchi} 26/10/2026, 0.2.0 (D32) \\
M2 — patch à chaud & remplacement du bloc d'écho en exécution, sans interruption ;
rollback démontré ; régression post-patch nulle & \textbf{franchi} 27/10/2026, 0.3.0 (D43) \\
M3 — entrées non fiables & un million de requêtes fuzzées sans crash ; proxy
transparent ; contrats renforcés & \textbf{franchi} 28/10/2026, 0.4.0 (D74) \\
M4 — tunnel chiffré & HTTPS en transit ; échec propre sur certificat invalide & \textbf{franchi} 29/10/2026, 0.5.0 (D99 — tunnel chiffré établi et certificat rejeté sur machine réelle, terrain.17-D ; balise scellée, arbre \code{m4/I} signé) \\
M5 — outil complet & interception d'un flux navigateur en VM isolée ; patch à chaud du
journal en production ; empreintes par build actives ; programme d'étapes et critères
dans \code{PLAN\_MILESTONES.md} & \textbf{franchi} 04/11/2026, \code{0.6.0} (Entrée 143, D148 : critère (b) de D136 vert sur pièces au run e4 ; le chantier (c) refermé au M6 sur le profil Linux) \\
M6 — portabilité inter-OS & suite de conformité verte ; \textbf{diff nul sur la couche
I} ; parité inter-profils vérifiée ; même oracle vert sur les deux profils & \textbf{franchi} 05/11/2026, \code{0.7.0} (Entrée 149, D154 : run terrain WSL2 vert sur pièces, recoup labo 22/22 — le noyau de la machine du développeur distinct de celui du labo) \\
M7 — portabilité inter-ISA & round-trip 100 \% sur la nouvelle ISA ; suite de conformité
verte ; patch à chaud conforme (scratch + \code{fence.i}) ; verrou multi-profils ; export
multi-profils vérifié & \textbf{franchi} 05/11/2026, \code{0.8.0} (Entrée 155 : run terrain M7-(d) vert sur pièces, recoup labo 29/29, export multi-profils de bout en bout — l'universalité démontrée : win \code{0.6.0} / linux \code{0.7.0} / riscv \code{0.8.0}) \\
M8 — syscalls directs (domaine \code{sys}) & la charte à huit règles et les spécifications sys-01…03 écrites-avant-code ; le harnais sur fixtures \code{ntdll} synthétiques (le labo ne connaît jamais le SSN du terrain) ; le motif du stub vérifié avant toute extraction ; le run terrain sur le VRAI \code{ntdll} — la même observable par les DEUX chemins ; la cérémonie de l'ère avant le build & \textbf{franchi} 05/11/2026, \code{0.9.0} (Entrées 156--164, D160--D165 : harnais 61 verts ; cérémonie de l'ère passée 12/12 ; run terrain m8d vert — 35 verts / 0 échec, la twin octet pour octet, le refus $-4$ du stub crocheté ; recoup labo 30/30 ; finale scellée sur recoup exigé comme pièce) \\
M9 — TLS 1.3 pur (domaine \code{tls13}) & la charte à huit règles et les 28 spécifications en six familles écrites-avant-code ; les primitives aux vecteurs officiels (FIPS 180-4, RFC 2104/4231, RFC 5869, RFC 7748, SP 800-38D) ; la poignée rejouée sur le transcript RFC 8448 ; le socle socket par syscalls directs (sys-04+) ; la cérémonie de l'ère avant \code{construire()} ; la session réelle — la ligne applicative chiffrée & \textbf{franchi} 07/11/2026, \code{1.0.0} (Entrées 170--194, D166--D170 : les 28 blocs notés-mesurés-prouvés aux vecteurs officiels, la poignée sur le transcript RFC 8448, le socle socket par syscalls directs ; le septième run terrain vert recoupé 34/34 sur pièces croisées, LA SESSION RÉELLE — la ligne applicative chiffrée échangée avec un serveur réel, le témoin « HTTP/1.1 200 OK » déchiffré par la machine, l'épinglage $-6$ nommé ; la finale scellée sur recoup exigé comme pièce, \code{proxy.exe} byte-identique \code{446279f3…}) \\
M10 — le test grandeur nature « Loupe » (domaine \code{loupe}) & coder une application réelle NEUVE entièrement avec TRIADE, de l'intention (la couche I) à l'octet livré (le kit terrain) ; les 18 spécifications écrites-avant-code avant tout octet de notation (R2) ; les héritages PAR LA CARTE byte-identiques aux adresses scellées (R7) ; l'oracle à scénarios vert — l'écran attendu octet par octet ; le kit déployé au coffre \code{/test} (modèle m9f) ; la balise \code{loupe 0.1.0} scellée sur run terrain recoupé & \textbf{ouvert} 07/11/2026 (D171) — la couche I complète (18 spécifications, D172) et les QUATRE chantiers de notation accomplis : le noyau (D173), le rendu + la navigation (D174), l'édition + la table de routage étendue de loupe-09 r2 (D175) — les 18 blocs NOTÉS-MESURÉS-PROUVÉS (59 + 53 + 45 contrôles verts aux trois graines) ; LA CÉRÉMONIE D'ÈRE du TGN PASSÉE SUR PIÈCES (le 14e arbre, la dette D166/D171 SOLDÉE — le garde 14/14) ; le kit terrain livré (M10-(d)) ; LE RUN (M10-(e)) REFUSÉ DEUX FOIS par le VRAI chargeur (0xc000007b : dev.1 sans BASERELOC, puis dev.2 — nécessaire mais pas suffisant) ; LA RETOUCHE dev.3 DÉPOSÉE : l'image ramenée DANS l'enveloppe prouvée de la lignée ET LE KIT PORTE SA CASCADE DIAGNOSTIQUE (le VRAI chargeur nommera l'axe au témoin) — le run rejoué attend le développeur \\
\bottomrule
\end{tabularx}
\end{table}

M0 est le jalon existentiel — il prouve la chaîne notation, octets, exécutable, exécution
sans compilateur. M2 est le jalon de vérité — le patch à chaud a motivé l'architecture.
M6 et M7 transforment la portabilité de promesse en propriété démontrée, en isolant chaque
variable ; la suite de conformité par profil en est désormais le critère d'acceptation
explicite. M8 et M9 sont les additions mesurées du programme d'après M7 (D160) : la
fondation d'abord — le domaine \code{sys} est la condition d'une pile TLS 1.3 qui
n'émettrait plus par l'IAT userland ; chaque famille s'appuie sur les résultats
mesurés de la précédente. M10 est l'épreuve de la MÉTHODE ENTIÈRE (D171) : non plus
une capacité, mais une application réelle neuve — l'outil du quotidien du projet —
développée de l'intention à l'octet livré, chaque geste prouvé ; la notation entière
(18 blocs) est notée-mesurée-prouvée, la cérémonie d'ère est PASSÉE (le garde à
quatorze arbres), le build est NÉ — et LE TERRAIN répond : deux refus nommés
(0xc000007b), la retouche dev.3 dans l'enveloppe prouvée de la lignée, et le kit
porte sa cascade diagnostic — le VRAI chargeur nomme l'axe fautif au témoin.

\paragraph{Le programme d'après M7 (D160).} La clôture v0.6.1 tenait les additions
mesurées « hors périmètre courant » ; la v0.7.0 lève la clause pour les deux
premières familles, sur résultats mesurés. \emph{M8, les syscalls directs} (domaine
\code{sys}) : la charte à huit règles — le SSN est une donnée, jamais un octet de
code ; le motif du stub est vérifié avant toute extraction ; un stub divergent est
un refus nommé — et la preuve d'observable : la MÊME ligne émise par les deux
chemins (l'IAT \code{kernel32} de hello-02 et le syscall direct \code{NtWriteFile}
de sys-03), comparées octet pour octet au fixture. \emph{M9, TLS 1.3 pur} (domaine
\code{tls13}, ère 1.0.0) : la grande œuvre de la souveraineté complète — la
cryptographie de session sans délégation, la couche I complète (28 spécifications,
six familles), la cérémonie d'ère passée (13 arbres, garde 13/13). L'ordre est la
fondation d'abord : une pile TLS pure qui émettrait encore par l'IAT userland
garderait la surface de hook qu'elle prétend éliminer.

\paragraph{Le test grandeur nature (D171).} Le protocole a prouvé chaque capacité,
M0 à M9 ; il reste à prouver qu'il DÉVELOPPE une application réelle NEUVE, de
l'intention à l'octet livré — le modèle RoomEq (v0.3) s'arrêtait au recodage d'une
application existante, le TGN exige la CRÉATION. L'application arbitrée est
« Loupe », l'éditeur hexadécimal natif au panneau d'empreintes sur le profil
primaire (\code{loupe.exe}, la lignée de balises propre \code{loupe 0.1.0-dev.1} au
premier build) : voir les octets, leurs lignes, leurs empreintes md5/sha-256, leur
place — l'outil du quotidien du projet. La découpe du plan est la loi : la cérémonie
d'ère précède le build (le miroir M8/M9 — Entrées 157/158), la notation s'accomplit
sous la dette nommée (D166/D171).

\paragraph{Archivé (D133, pour mémoire).} Le jalon de l'outil complet s'ouvrit sur trois
chantiers mesurables : (a) \emph{empreintes par build actives} — chaque build consigne
l'empreinte de son artefact dans le tableau de bord et le journal de travail, la
comparaison inter-build devient un contrôle automatique ; (b) \emph{interception d'un
flux navigateur en VM isolée} — le proxy transparent M3 porte le tunnel M4 au scénario
réel complet ; (c) \emph{patch à chaud du journal en production} — l'écriture vivante
du superviseur terrain sous patch, le retour arrière prêt. L'ordre est dicté par la
dépendance : (a) arme la preuve dont (b) et (c) ont besoin pour être auditables.

\section{Correspondance des versions}\label{sec:correspondance}

\subsection{Traçabilité des évolutions}\label{sec:evolutions}

\begin{table}[htbp]
\centering\small
\caption{Les dix-sept évolutions de la v0.4.0, tracées depuis la v0.3.}
\label{tab:correspondance}
\begin{tabularx}{\textwidth}{@{}l Y Y l@{}}
\toprule
\textbf{Évolution} & \textbf{En v0.3} & \textbf{En v0.4.0} & \textbf{Statut} \\
\midrule
Couche I normalisée & « contrats légers » sans gabarit & gabarit complet
(§\ref{sec:coucheI}) & renforcé (D19) \\
Taxonomie des défauts & absente & classes T1–T6 (§\ref{sec:taxonomie}) & nouveau (D19) \\
Threat model & 7 menaces & 13, dont 6 nouvelles & étendu (D20) \\
R10 & checksums au lancement & + signature, anti-rollback, conformité, fail-closed
explicite & renforcé (D20) \\
Optimisation prématurée & interdite implicitement & interdiction normative R11 &
explicité (D21) \\
Journal d'optimisation & pratique de démonstration & format OPT-NN & officialisé (D21) \\
Fragmentation & 4 fragments & + norme « aucune lecture totale » & consolidé (D24) \\
Localisation par intention & étape 2 implicite & balayage interdit, formalisé &
renforcé (D24) \\
Budgets tokens & 64 / 3 000 / 6 000 & localiser amendé après mesure M1, instrumentés & consolidé (D24, D31) \\
Suite de conformité profil & implicite & suite obligatoire (§\ref{sec:conformite}) &
nouveau (D22) \\
Parité inter-profils & critères dispersés & définition formelle (§\ref{sec:parite}) &
formalisé (D22) \\
Reverse-specs & absente (pratique démo) & officialisées (§\ref{sec:reversespecs}) &
officialisé (D23) \\
Pré-versions & absente (pratique \code{-dev.1}) & dev / rc / final & officialisé (D23) \\
Fixtures partagées & absente (pratique installateur) & contractuelles & officialisé
(D23) \\
Traçabilité I–C–B–tests & implicite & table normative + audit & formalisé (D24) \\
Journal des décisions & pratique & format normalisé & normalisé (D24) \\
Tableau de bord & indicateurs dispersés & métriques unifiées & nouveau (D24) \\
\bottomrule
\end{tabularx}
\end{table}

\subsection{Modifications explicitement rejetées}\label{sec:rejets}

Conformément au cahier de révision, les modifications suivantes ont été examinées et
\textbf{rejetées, donc non appliquées} :

\begin{table}[htbp]
\centering\small
\caption{Modifications examinées et rejetées.}
\label{tab:rejets}
\begin{tabularx}{\textwidth}{@{}Y Y@{}}
\toprule
\textbf{Modification} & \textbf{Motif du rejet} \\
\midrule
Modifier la hiérarchie Sécurité $>$ Performance $>$ Optimisation $>$ Coût en tokens &
invariant de révision (D19) — condition même de l'arbitrage \\
Introduire une couche d'abstraction à l'exécution (bytecode, VM) & trahit la fonction
objectif Performance ; la parité s'obtient par vérification, pas par couche \\
Promettre la portabilité universelle immédiate & la portabilité se vérifie par profil
(M6/M7), elle ne se déclare pas \\
Introduire des mécanismes de furtivité face aux EDR & hors périmètre assumé —
transparence totale de l'outil \\
Affaiblir R5, R8, R9 ou R10 & la révision ne contient que des renforcements ; R1–R9
repris à l'identique \\
Anti-rollback par contrôle en ligne au chargement & R10 doit échouer proprement hors
ligne — aucune dépendance réseau \\
Optimisation validée sans mesure & R11 : hot path + mesure avant/après + contrat stable \\
\bottomrule
\end{tabularx}
\end{table}


\subsection{Correspondance v0.4.0 $\to$ v0.5.0}\label{sec:evolutions050}

\begin{table}[htbp]
\centering\small
\caption{Les huit évolutions de la v0.5.0, tracées depuis la v0.4.0.}
\label{tab:correspondance050}
\begin{tabularx}{\textwidth}{@{}l Y Y l@{}}
\toprule
\textbf{Évolution} & \textbf{En v0.4.0} & \textbf{En v0.5.0} & \textbf{Statut} \\
\midrule
Validation terrain & constats dispersés au journal des décisions & couche normée :
boucle d'évidence, journal d'exécution, sondes contractualisées, diff chirurgical
(§\ref{sec:terrain}) & nouveau (D43) \\
Taxonomie des défauts & classes T1--T6 & + classes X1--X8, invisibles à l'émulation
(§\ref{sec:taxonomieX}) & étendu (D35--D42) \\
État des jalons & M0 planifié & M0, M1, M2 franchis — preuves terrain archivées
(§\ref{sec:jalons}) & constaté (D26, D32, D43) \\
Budgets de tokens & 40 / 3 000 / 6 000 (cibles) & 64 / 3 000 / 6 000 — corrigé et mesuré :
51 / 719 / 1 088 & corrigé (D31) \\
Suite de conformité & définie (7 éléments) & exécutée à chaque scellage rc/final — 124/124
sur le profil primaire & mesuré (D39--D43) \\
Volume d'oracle & 12 000 cas par bloc critique (référence) & campagnes effectives
consignées : 76 010 (M1), 50 533 (M2), 86 (démo) & consolidé \\
Preuves terrain & non normées & archives brutes (journaux octet-exacts, empreintes)
référencées au journal des décisions & officialisé (D43) \\
Fuzzing & mentionné (menace « fuzzing boîte noire ») & jalon M3 ouvert : fuzzing à
propriétés sur la surface d'entrées non fiables & structuré (D45) \\
\bottomrule
\end{tabularx}
\end{table}

\paragraph{Modifications explicitement rejetées (v0.5.0).} Ajouter une règle R12
« validation terrain obligatoire » — rejetée : la boucle d'évidence découle de R5 et des
critères de jalon, une règle de plus diluerait le corpus sans rien exiger de nouveau.
Étendre le fuzzing à tous les blocs scellés — rejetée : M3 cible la surface d'entrées non
fiables, les blocs à appelants de confiance relèvent des vecteurs existants. Reconstruire
les artefacts scellés 0.2.0 / 0.3.0 pour la documentation — rejetée : les balises finales
sont immuables, la documentation cite les empreintes. Fusionner les classes X dans les
classes T — rejetée : les deux familles appellent des remèdes distincts, écrire le contrat
(T) contre enrichir le modèle d'environnement et falsifier (X).

\subsection{Correspondance v0.5.0 $\to$ v0.5.1}\label{sec:evolutions051}

\begin{table}[htbp]
\centering\small
\caption{Les évolutions de la v0.5.1, tracées depuis la v0.5.0.}
\label{tab:correspondance051}
\begin{tabularx}{\textwidth}{@{}l Y Y l@{}}
\toprule
\textbf{Évolution} & \textbf{En v0.5.0} & \textbf{En v0.5.1} & \textbf{Statut} \\
\midrule
Phrases maîtresses & instruction de cérémonie seule (D27/D30) — quatre
réinitialisations du poste ont montré la limite & coffre-fort TRIADE-CF1 : conteneur
chiffré par projet, export/import, restauration post-réinitialisation & nouveau (D75) \\
État des jalons & M0–M2 constatés, M3 ouvert & M3 franchi 28/10/2026 — campagne
million (D72), conformité permanente (D73), scellage 0.4.0 (D74) & constaté (D72–D74) \\
Modèle de menace & 13 menaces, dont le vol de l'export & + vol du coffre-fort des
phrases — sans la phrase de contrôle, indéchiffrable & étendu (D75) \\
\bottomrule
\end{tabularx}
\end{table}

\paragraph{Modifications explicitement rejetées (v0.5.1).} Conserver la seule
instruction de coffre développeur — rejeté : le statu quo que les quatre réinitialisations
ont démenti. Stocker les phrases en clair dans \code{.keys/} — rejeté : l'état
opérationnel du poste n'est pas un coffre. Embarquer la graine Ed25519 dans le coffre —
rejeté : contraire à D30, jamais confiée. Confier le conteneur au seul gestionnaire de
mots de passe du développeur sans copie de coffre projet — rejeté : la continuité exige
les deux copies, le conteneur étant indéchiffrable sans la phrase.

\subsection{Correspondance v0.5.1 $\to$ v0.5.2}\label{sec:evolutions052}

\begin{table}[htbp]
\centering\small
\caption{Les évolutions de la v0.5.2, tracées depuis la v0.5.1.}
\label{tab:correspondance052}
\begin{tabularx}{\textwidth}{@{}l Y Y l@{}}
\toprule
\textbf{Évolution} & \textbf{En v0.5.1} & \textbf{En v0.5.2} & \textbf{Statut} \\
\midrule
Garde de la phrase de contrôle & détenue par le seul développeur, jamais stockée,
affichée une seule fois (D75) & sauvegardée durablement dès production :
\code{PHRASE-CONTROLE-\allowbreak COFFRE.txt} au dépôt (git, GitHub, miroir Drive) et au racine du
coffre Drive, AVANT toute autre action ; le fichier est la remise & modifié (D81) \\
Menace « vol du coffre-fort » & sans la phrase, indéchiffrable & la phrase vit sur le
coffre Drive partagé par lien : confidentialité au repos ramenée au contrôle d'accès au
dossier ; une phrase commise en git y demeure à vie, la rotation reste le remède &
étendu (D81) \\
État des jalons & M0–M3 franchis, M4 à venir & M4 ouvert (D77 — couche I
écrite-avant-code : huit spécifications tls, deux révisions http) ; M0–M3 inchangés &
constaté (D77) \\
\bottomrule
\end{tabularx}
\end{table}

\paragraph{Modifications explicitement rejetées (v0.5.2).} Conserver la remise par
affichage unique — rejeté : le statu quo que la session morte de D80 a démenti (la
phrase ne vivait plus que dans le bac à sable). Sceller la phrase de contrôle sous une
phrase antérieure — rejeté : déplace le problème, la phrase de secours exigerait sa
propre garde. Exporter la seule copie chiffrée du coffre sans garder la phrase — rejeté :
le conteneur indéchiffrable protège de l'exposition, pas de la perte, or c'est la perte
qui est le risque vécu. Restreindre d'office le partage du dossier Drive — rejeté : la
décision appartient au développeur, hors du champ de l'agent (consignée en
recommandation).

\subsection{Correspondance v0.5.2 $\to$ v0.6.0}\label{sec:evolutions060}

\begin{table}[htbp]
\centering\small
\caption{Les évolutions de la v0.6.0, tracées depuis la v0.5.2.}
\label{tab:correspondance060}
\begin{tabularx}{\textwidth}{@{}l Y Y l@{}}
\toprule
\textbf{Évolution} & \textbf{En v0.5.2} & \textbf{En v0.6.0} & \textbf{Statut} \\
\midrule
Coffre-fort & conteneur TRIADE-CF1 à phrase (D75), garde durable de la phrase (D81) &
chaîne complète : CF2 (clé maîtresse Shamir, cinq gardes seuil 3, sentinelle/porte
D120) puis CF3 (serrure à triple scellement, piège sur média, cycle en deux actes,
sentinelle instantanée par invocation) — le coffre réel a migré par cérémonies
(30--31/10/2026) ; un coffre neuf naît CF3 & ajouté (D86, D106, D121, D127, D133) \\
Souveraineté de l'intention & exigence consignée (D87, programme ouvert) & le sceau de
l'Intention vit : dix arbres signés par la clé Ed25519 hors-ligne, cérémonie hors-ligne
autoportée, gardes de build et de patch, re-cérémonie réelle accomplie (D126) &
ajouté (D87, D105, D126) \\
Cérémonies et kit & non couverts (l'outil \code{coffre\_fort.py} seul) & maître de
cérémonie graphique (relais, zéro ligne de commande), trio D122 (GUI + boîte noire +
rapport), kit Drive identique au miroir par construction (D130), lanceurs Windows/POSIX &
ajouté (D107--D132) \\
Attestation TEE (axe 1) & exigence consignée (D85) & inchangée : exigence consignée —
l'infrastructure sans racine de confiance matérielle ne peut pas la satisfaire et le dit &
constaté (D133) \\
Menaces & vol du coffre-fort, rollback binaire & lignes nouvelles : expropriation, evil
maid sur conteneur/média, roll-back du couple conteneur + piège & étendu (D133) \\
Jalons & M4 ouvert (D77) & M4 franchi (0.5.0 scellée, terrain.17-D) ; M5 ouvert (D133) ;
axes 2-3-4 du programme v0.6.0 prouvés terrain & constaté \\
\bottomrule
\end{tabularx}
\end{table}

\paragraph{Modifications explicitement rejetées (v0.6.0).} Supprimer la sentinelle en
 guise de simplification — rejeté : l'exigence D120 (« indispensable au déverrouillage »)
 est réalisée \emph{par invocation} (la sentinelle instantanée est la serrure, impossible
 à oublier de lancer) ; ce qui meurt est le démon résident, usurpable par faux
 heart-beat. N'imposer la chaîne de migration aux coffres neufs — rejeté : un coffre neuf
 naît CF3 (\code{init}, époque 1) ; la chaîne CF1 $\to$ CF2 $\to$ CF3 est l'histoire du
 coffre réel, pas un rite. Signer la présente spécification avant passation — reporté :
 les sources \code{docs/} sont hors périmètre du sceau courant ; elles seront signées à
 leur propre version (modèle par-version, \code{v06/I/README.md} § 1). Passer à v0.6.0 en
 attendant l'axe 1 — arbitrée par le développeur : la passation se fait avec l'axe 1
 consigné (D133), le plan le déclarait mûre dès l'écriture des trois composants.

\subsection{Correspondance v0.6.0 $\to$ v0.6.1}\label{sec:evolutions061}

\begin{table}[htbp]
\centering\small
\caption{Les évolutions de la v0.6.1, tracées depuis la v0.6.0 (D137).}
\label{tab:correspondance061}
\begin{tabularx}{\textwidth}{@{}l Y Y l@{}}
\toprule
\textbf{Évolution} & \textbf{En v0.6.0} & \textbf{En v0.6.1} & \textbf{Statut} \\
\midrule
Niveau de naissance d'un coffre neuf & un coffre neuf naît CF3 — \code{init} direct à
l'époque 1 (D133) & le niveau (CF1/CF2/CF3) est arbitré par le développeur à la création
 de son projet : les trois naissances, le guide et les obligations par niveau
 (§\ref{sec:initniveau}) ; la chaîne de migration demeure disponible à tout moment &
révisé (D137) \\
Le coffre réel & CF3 (époque 1, checksum \code{0cc89a96}) & inchangé — la révision est un
droit de création, pas un geste sur le coffre vivant & constaté \\
Chaîne PDF et Livre du Coffre & spécification v0.6.0 ; Livre v1.0 & chaîne spec\_v061
(l'erratum typographique « tbp] » page 31 réparé au passage) ; Livre v1.1 — runbook
« coffre neuf » et FAQ terrain alignés sur D137 & révisé (D137) \\
Le reste de la spécification & — & inchangé : aucune autre règle, aucune couche
d'exécution ajoutée, hiérarchie inchangée & constaté \\
\bottomrule
\end{tabularx}
\end{table}

\paragraph{Modifications explicitement rejetées (v0.6.1).} Imposer CF1 par défaut en
guise de souplesse — rejeté : remplacer une imposition par une autre ne rend pas la
décision au développeur. Rendre le niveau mutable sans cérémonie — rejeté : le niveau
change par MIGRATION uniquement, acte explicite et fail-closed (l'anti-retour d'époque
et le verrou croisé demeurent). Promouvoir l'\code{init} de CF2 (laboratoire) en
cérémonie réelle dans cette révision — rejeté : la spécification scellée CF2 fait foi ;
la naissance CF2 passe par \code{init} CF1 puis \code{migrer}, tout est prouvé ;
l'arbitrage est consigné si le besoin naît. Toucher au code des trois outils — rejeté :
les trois actes de création existent et sont prouvés ; c'est le droit de choisir qui
manquait, pas la mécanique.

\subsection{Correspondance v0.6.1 $\to$ v0.7.0}\label{sec:evolutions070}

\begin{table}[htbp]
\centering\small
\caption{Les évolutions de la v0.7.0, tracées depuis la v0.6.1 (D169).}
\label{tab:correspondance070}
\begin{tabularx}{\textwidth}{@{}l Y Y l@{}}
\toprule
\textbf{Évolution} & \textbf{En v0.6.1} & \textbf{En v0.7.0} & \textbf{Statut} \\
\midrule
Programme d'après M7 & clause de clôture : syscalls directs et TLS 1.3 pur « hors
périmètre courant » & le programme D160 entre en spécification : la clause est
LEVÉE NOMMÉMENT pour les deux premières familles — M8 (domaine \code{sys})
franchi, M9 (domaine \code{tls13}) ouvert ; les profils additionnels demeurent
hors périmètre & révisé (D160, D166) \\
Sceau de l'Intention & dix arbres sous le sceau & treize arbres (m0--m8, tls13,
demo, terrain, v06) ; les cérémonies d'ère entrent au protocole : la dette nommée
à la naissance (D166 règle 8), la charge préparée au geste qui sert la main
(D168), la continuité de clé prouvée par le déterminisme Ed25519 (D163) & étendu
(D162, D166, D168) \\
État des jalons & M0--M6 franchis, M7 ouvert & M0--M8 franchis (balises 0.2.0 à
0.9.0), M9 ouvert — la cérémonie de l'ère M9 passée (garde 13/13) & constaté
(Entrées 155--168) \\
Chaîne PDF & spec\_v061 & chaîne spec\_v070 ; l'archive v0.6.1 demeure (dernier
build consigné : sha256 \code{4aac10f1…}, 984 555 o) & révisé (D169) \\
\bottomrule
\end{tabularx}
\end{table}

\paragraph{Modifications explicitement rejetées (v0.7.0).} Réviser la spécification
à chaque franchissement de jalon — rejeté : le versionnement suit les RÉVISIONS,
jamais les franchissements (l'état des jalons vit aux sources de travail,
régénérées à chaque geste ; la v0.7.0 consolide un PROGRAMME, pas un simple
constat). Signer les sources \code{docs/} sous le sceau courant — rejeté : le
modèle par-version demeure (\code{v06/I/README.md} § 1). Ouvrir les profils
additionnels (\code{aarch64-macos-macho}, \code{wasm32-wasi}) dans la même
révision — rejeté : chacun répétera M6/M7 sur résultats mesurés, hors périmètre.
Fusionner les cérémonies d'ère dans la cérémonie CF3 — rejeté : trois gardes,
trois menaces ; le sceau garde l'intention, le coffre garde les secrets.

\subsection{Correspondance v0.7.0 $\to$ v0.8.0}\label{sec:evolutions080}

\begin{table}[htbp]
\centering\small
\caption{Les évolutions de la v0.8.0, tracées depuis la v0.7.0 (D176).}
\label{tab:correspondance080}
\begin{tabularx}{\textwidth}{@{}l Y Y l@{}}
\toprule
\textbf{Évolution} & \textbf{En v0.7.0} & \textbf{En v0.8.0} & \textbf{Statut} \\
\midrule
Programme d'après M9 & M9 ouvert : couche I complète, la cérémonie passée, les
chantiers (b) à (f) attendent & M9 FRANCHI : les 28 blocs notés-mesurés-prouvés,
le septième run terrain vert recoupé, LA SESSION RÉELLE (la ligne applicative
chiffrée avec un serveur réel), la finale \code{1.0.0} scellée sur recoup exigé
comme pièce (D170) & révisé (D170) \\
Le test grandeur nature & — & LE TGN ouvert (D171) : « Loupe », l'éditeur
hexadécimal au panneau d'empreintes — coder une application réelle NEUVE
entièrement avec TRIADE ; le domaine \code{loupe} fondé, la couche I de 18
spécifications posée (D172), les QUATRE chantiers de notation accomplis
(D173/D174/D175), la cérémonie d'ère PASSÉE (le garde 14/14), le build
\code{loupe.exe 0.1.0-dev.1} NÉ ; LE RUN terrain REFUSÉ DEUX FOIS (0xc000007b —
dev.1 sans BASERELOC, dev.2 : nécessaire mais pas suffisant), la retouche dev.3
DANS l'enveloppe prouvée de la lignée, LE KIT PORTE SA CASCADE DIAGNOSTIQUE & mis
à jour (D171--D175 + la révision de terrain v0.8.1) \\
Sceau de l'Intention & treize arbres, les cérémonies d'ère M8/M9 passées & treize
arbres + LA DETTE DU QUATORZIÈME nommée (l'ère du TGN : l'arbre \code{loupe} et
le registre figé étendu attendent la clé hors-ligne AVANT tout build —
§\ref{sec:ceremoniesere}) ; le garde demeure à 13/13 jusque-là & étendu
(D166/D171/D175) \\
Projet fils COFFRE-FORT & développement différé (D135) & DÉGELÉ en parallèle du
TGN (D171) : le fils vit son rythme, la frontière fondatrice intacte — aucun code
TRIADE forké, aucune dépendance vivante & révisé (D171) \\
Chaîne PDF & spec\_v070 & chaîne spec\_v081 (née spec\_v080 — la révision de
terrain v0.8.1) ; les archives v0.7.0 et v0.8.0 demeurent (derniers builds
consignés : \code{dd36eb41…} et \code{95226c2b…}) & révisé (D176, v0.8.1) \\
\bottomrule
\end{tabularx}
\end{table}

\paragraph{Modifications explicitement rejetées (v0.8.0).} Réviser la spécification
à chaque franchissement de jalon — rejeté : le versionnement suit les RÉVISIONS (la
v0.8.0 consolide le PROGRAMME d'après M9 : M9 franchi ET le TGN ouvert, deux
moteurs d'une même passation). Mettre l'application « Loupe » sous le sceau AVANT
la cérémonie — rejeté : le garde suit le calendrier de la cérémonie, jamais
l'inverse (§\ref{sec:ceremoniesere}). Ajouter un profil de cible pour le TGN —
rejeté : le profil primaire \code{x86-64-win11-pe} est l'épreuve ; l'universalité
est déjà prouvée (M6/M7). Sceller la 0.1.0 de Loupe au build — rejeté : la balise
de développement naît au build (geste de BUILD jamais scellage, miroir
M6/M7/M8) ; la 0.1.0 se scellera sur le run terrain vert recoupé.

\subsection{Correspondance v0.8.0 $\to$ v0.8.1}\label{sec:evolutions081}

La v0.8.1 est UNE RÉVISION DE TERRAIN : aucune règle modifiée, aucune couche
d'exécution ajoutée, hiérarchie inchangée — le protocole CONSIGNE ce que le
grandeur nature a dit entre la livraison du kit et le run rejoué (le terrain est
une lecture d'exécution : son enseignement se consigne, §\ref{sec:validationterrain}).

\begin{table}[tbp]
\centering
\small
\begin{tabularx}{\textwidth}{>{\raggedright\arraybackslash}p{3.4cm} X}
\textbf{Évolution} & \textbf{Le fait consigné (v0.8.1)} \\
\midrule
La cérémonie d'ère du TGN & PASSÉE SUR PIÈCES (09/11, Entrée 201) : le garde à
QUATORZE arbres verts, la dette D166/D171 SOLDÉE, la remise du développeur
INSTALLÉE au coffre \code{/test} — le registre figé étendu porte l'arbre
\code{loupe} \\
Le build du domaine & \code{loupe.exe 0.1.0-dev.1} NÉ (Entrée 202) : la
composition intégrale comptée, la carte génération 13, le verrou croisé 26/26,
les trois harnais rejoués verts au geste du build (157 contrôles) \\
LE REFUS DU TERRAIN, DEUX FOIS & le run M10-(e) a REFUSÉ dev.1 puis dev.2 du
même refus nommé 0xc000007b au geste de cartographie (LoadLibraryEx
DONT\_RESOLVE — le VRAI chargeur Windows) : dev.1 portait l'oubli du BASERELOC
minimal (l'Entrée 204 : la table rendue au plan, le miroir M1..M9) ; dev.2
PROUVE que la table était NÉCESSAIRE mais PAS SUFFISANTE — l'image sortait de
l'enveloppe des HUIT artefacts de la lignée passés au terrain : le trou de page
0x3000--0x4000, le BSS massif \code{.data} (VS $\gg$ RS), le sous-système GUI \\
La retouche dev.3 & l'image ramenée DANS l'enveloppe prouvée : \code{.rdata}
étendue 0x1000 $\to$ 0x2000 (plus AUCUN trou de page), \code{.data} VS~$=$~RS
(le RAW entier, plus AUCUN BSS) — AUCUNE adresse scellée ne bouge, AUCUN octet
de notation (R2) ; le build rejoué : garde 14/14, oracle 43 verrous, harnais
59+53+45 verts, \code{loupe.exe 0.1.0-dev.3} : 127 488 o, sha256
\code{b494ff90…} \\
LA CASCADE DIAGNOSTIQUE & le kit dev.3 PORTE SON BISSECTEUR : le scénario
reconstitue UN À UN les axes de dev.2 sur des copies éphémères \%TEMP\%
(l'artefact n'est jamais écrit) — P0 le contrôle, P1 le trou, P2 le BSS, P3 le
sous-système, P4 la géométrie entière — chaque verdict VERTE/REFUS consigné au
témoin avec GetLastError + NTSTATUS ; ET le mode d'erreur du processus coupe
les boîtes bloquantes du chargeur : le refus est nommé PAR LE TÉMOIN, jamais
muré derrière un bouton OK \\
L'enseignement consigné & le TGN fait exactement ce pour quoi il existe :
l'application vit chez son utilisateur, et SEUL le VRAI chargeur pouvait nommer
le défaut (le labo est Linux : le harnais PARSE l'image, il ne la charge
jamais) — l'enveloppe prouvée de la lignée devient une NORME DE FAIT pour
l'éditeur PE du build : ce que le terrain n'a pas prouvé, le build ne l'ose pas \\
\bottomrule
\end{tabularx}
\caption{La révision de terrain v0.8.1 : les Entrées 201--205 du journal de
travail, consignées au protocole. Aucune règle modifiée.}
\end{table}

\paragraph{Modifications explicitement rejetées (v0.8.1).} Blâmer le terrain —
rejeté : le refus nommé est un ENSEIGNEMENT, la preuve vivante que seule
l'exécution réelle tranche (l'enveloppe prouvée devient la norme de fait).
Étendre l'éditeur PE au-delà de la lignée sur spéculation — rejeté : le TGN ne
négocie pas avec le chargeur ; ce que les HUIT n'ont pas prouvé, la retouche ne
l'ose pas (la sonde tranchera si l'axe devait être légitime). Reporter le run au
labo — rejeté : Windows seul peut nommer le refus du chargeur Windows ; le kit
repart chez l'utilisateur avec son bissecteur.

\section{Journal des décisions}\label{sec:decisions}

Le journal complet vit dans \code{docs/DECISIONS.md} du dépôt du projet. La
v0.8.0 consolide le programme d'après M9 : D170 (le jalon M9 franchi — la finale
\code{1.0.0} scellée sur le recoup exigé comme pièce, la session réelle la ligne
applicative chiffrée), D171 (le test grandeur nature ouvert — « Loupe » arbitrée,
le domaine \code{loupe} fondé, la dette du quatorzième arbre nommée à la
naissance, le projet fils COFFRE-FORT dégelé), D172 (les 18 spécifications avant
tout octet de notation), D173/D174/D175 (les chantiers de notation — noyau,
rendu + navigation, édition : les 18 blocs notés-mesurés-prouvés) et D176 (la
présente consolidation). Synthèse des débuts :
D1–D10 (v0.2) fixent l'oracle éphémère, la notation, les contrats légers, la
cryptographie, le verrou croisé, les variantes par build, l'interface système et la
granularité ; D11–D13 (infrastructure) la sauvegarde triple et systématique ; D14–D16
(v0.3) l'universalité par profils, la cryptographie paramétrable et l'export de
continuité ; D17 la régénération systématique du PDF formel ; D18 l'édition académique
(présentation IMRAD du noyau v0.3, sans valeur normative) ; D19–D24 (v0.4.0, format
normalisé) la présente révision :

\begin{itemize}[label={}, leftmargin=1.2em]
\item \textbf{D19} : couche I normalisée, taxonomie T1–T6, hiérarchie de la fonction
objectif verrouillée comme invariant de révision ;
\item \textbf{D20} : menace étendue (binaire altéré, carte corrompue, rollback,
agent défaillant, données embarquées, environnement) et R10 renforcée (checksums,
signature, anti-rollback, fail-closed explicite, vérification locale) ;
\item \textbf{D21} : R11 optimisation mesurée et journal d'optimisation OPT-NN ;
\item \textbf{D22} : suite de conformité par profil et parité inter-profils à contrat
constant ;
\item \textbf{D23} : officialisation des reverse-specs, des pré-versions dev / rc /
final et des fixtures partagées ; numérotation à trois segments ;
\item \textbf{D24} : traçabilité I $\leftrightarrow$ C $\leftrightarrow$ B
$\leftrightarrow$ tests, journal des décisions normalisé, tableau de bord de conformité.
\end{itemize}

Deux décisions postérieures à la révision complètent le tableau, sans valeur
normative : D25 confie la production du présent PDF formel à la chaîne
LaTeX/Tectonic (générateur ReportLab retiré, archives des versions antérieures
conservées) ; D26 consigne le jalon M0, première étape de la feuille de route,
menée à terme le 26 octobre 2026 (chaîne artefact implémentée, \code{hello.exe}
exécuté sur Windows 11, carte v0 signée, quatre critères de sortie acquis).

Les décisions du jalon M1 poursuivent le journal : D27 (continuité des clés après
réinitialisation du poste : régénération et centralisation, preuve par artefact
byte-identique), D28 (plan d'adressage par zones et carte projet multi-artefacts),
D29 (substrat d'émulation Unicorn pour l'oracle différentiel local), D30 (troisième
réinitialisation : D27 réappliquée, premier export de continuité généré) et D31
(budget de localisation amendé à 64 tokens : la mesure M1 a constaté que le gabarit
canonique de la présente révision mesurait déjà 43 tokens en cl100k\_base —
l'instrumentation a fait son travail de révélateur). D32 clôt le jalon : la
validation terrain de \code{selftest.exe} sur Windows 11 est constatée (ligne unique
du contrat libc-08 rendue à l'identique, 21 vecteurs sur 21), la balise devient
\code{0.2.0} finale — immuable — et le contrôle anti-rollback de la reconstruction
est étendu aux balises finales ; le jalon de vérité M2 (patch à chaud) peut
commencer.

La campagne M2 alimente ensuite le journal au rythme du terrain : D33 (base préférée
par ancre RIP-relative — image indépendante de son adresse de chargement), D34
(échelle d’auto-hébergement : la vérification migre vers le natif, chantier m7), D35
(sémantique du chargeur — FirstThunk : la première classe X de la taxonomie du
§\ref{sec:taxonomieX}), D36 à D38 (les trois premières itérations de la démonstration
graphique : lecture de structure, canal de vérité, ABI caller-saved — classes X2 à X4),
D39 à D42 (la cascade : protections mémoire, écart émulation/terrain, zone d’accueil
ABI, état hérité — classes X5 à X8), D43 (M2 clos : boucle § 5 démontrée sur machine
réelle, retrait des sondes, scellage final 0.3.0 — la validation terrain devient
couche normée). D44 consigne la présente révision v0.5.0 ; D45 ouvre le jalon M3
(contrats renforcés sur entrées non fiables, fuzzing à propriétés).

La campagne M3 poursuit le journal jusqu'à la clôture du jalon : D46 (couche I écrite
avant le code — treize spécifications du domaine http), D47 (proxy.exe scellé),
D48–D50 (lignée terrain : le superviseur un-exe et sa chaîne de lecture), D51–D71 (la
chasse aux tueurs de runs en seize générations — ABI shadow space, pile du proxy,
réponses Google, tampon de composition — jusqu'au T3 démontré de bout en bout sur
machine réelle), D72 (campagne million : 10\^6 requêtes fuzzées sans crash, reproductible),
D73 (le fuzzing rejoint la suite de conformité — sentinelle perpétuelle), D74 (M3 clos :
cristallisation 38/38 $\times$ 10/10, scellage de la balise finale 0.4.0, lignée de clés
régénérée avec preuve byte-identique). D75 consigne la révision v0.5.1 : le
coffre-fort des phrases maîtresses — conteneur chiffré TRIADE-CF1 centralisant toutes
les phrases par projet, sur la pile cryptographique de la carte au repos, avec autotest
de falsification et restauration post-réinitialisation.

La suite du journal couvre l'hygiène et l'ouverture de M4 : D76 (nettoyage du coffre
Drive — douze livrables périmés archivés, miroir rendu fidèle à git, incident du reliquat
consigné), D77 (M4 ouvert par sa couche I écrite-avant-code — huit spécifications
tls et deux révisions http, chartre des contrats cryptographiques), D78 (le reliquat
re-qualifié : il portait la phrase COURANTE, exposée environ 64 minutes — recommandation
de rotation), D79 (M4 étape (b) engagée — trois feuilles pures notées, mesurées et
prouvées par différentiel de spécification et falsification), D80 (rotation de lignée
exécutée — la phrase exposée retirée du service, carte m3 re-scellée à faits inchangés,
coffre-fort roté). D81 consigne la présente révision v0.5.2 : la phrase de contrôle du
coffre-fort passe sous garde durable — sauvegardée au dépôt et au coffre Drive AVANT
toute autre action dès sa production, le fichier étant la remise.

La série de souveraineté ouvre la génération suivante : D84 place le dépôt sous licence
BSL 1.1 et nomme la menace d'expropriation ; D85 (attestation de non-exfiltration), D86
(interrupteur d'homme mort — Shamir, sentinelle, destruction des fragments) et D87
(filigrane de l'Intention — clé Ed25519 hors-ligne) fondent le programme v0.6.0. Puis
vient l'ère des livraisons : D105 (axe 3 livré en laboratoire — spec écrite-avant-code,
cérémonie hors-ligne autoportière, dix arbres, falsifiabilité 47/47) ; D106 (axe 2 livré
— coffre CF2, Shamir sur générateur primitif 0x03, sentinelle inotify + /proc à double
seuil, destruction trois passes + PV, migration CF1 $\to$ CF2) ; D107 (le maître de
cérémonie graphique — un fichier, deux profils, zéro ligne de commande) ; D108 (le kit
de cérémonie au Drive) ; D110 (la phrase de contrôle appartient au développeur :
alphanumérique, sans plafond) ; D111--D118 (l'attachement nommable, les trois chemins,
la boîte noire du GUI, le verdict ancré, le manifeste au contenu canonique LF, le fil
de cérémonie par preuves, le scénario avant tout, la page unique et les lanceurs).

Le terrain fait ensuite la loi : D119 (le dossier de cérémonie naît avec la clé ; fsync
POSIX-seul) ; D120 (la sentinelle est exigée au déverrouillage — la porte ferme le
coffre quand le témoin meurt) ; D121 (la serrure à triple scellement et le cycle en deux
actes : CF3 remplace le démon par un piège hors machine — les deux demandes du terrain
servies ensemble) ; D122 (le kit unique et le trio systématique : une GUI, une boîte
noire, un bouton « Envoyer le rapport »). L'exécution réelle suit : phrase changée
(D110), coffre réel migré CF1 $\to$ CF2 (31/10 — cinq gardes, seuil 3, trois copies
hors-ligne), re-cérémonie du sceau accomplie par le développeur depuis son clone
(D126 — commit 20f40bc, dix arbres verts), outillage de migration CF2 $\to$ CF3 livré
(D127 — \code{migrer}, \code{detruire-cf2}, GUI de migration, lanceurs), la chaîne
réelle achevée (30--31/10 : époque 1, checksum \code{0cc89a96}, verdict à froid VERT
deux fois). Les correctifs d'habit guidés par le terrain ferment la série : D129 (le
pilier droit), D130 (le kit ne dérive plus du miroir), D131 (un bouton d'attachement
par page), D132 (la grille du GUI de migration — chaque bouton « Parcourir » sur la
ligne de son champ).

\textbf{D133} consigne la révision v0.6.0 : la passation du programme de
souveraineté cryptographique — arbitrée par le développeur, les trois composants étant
écrits et la chaîne réelle prouvée sur terrain ; l'axe 1 (D85) demeure une exigence
consignée. D133 clarifie en outre la position normative de la sentinelle : l'exigence
D120 est réalisée \emph{par invocation} dans CF3 (la sentinelle instantanée est la
serrure — aucun démon à lancer, aucun moteur de plateforme à attendre), et un coffre
neuf naît CF3 par \code{init} direct : la chaîne de migration appartient à l'histoire du
coffre réel. Le jalon M5 est ouvert dans le même geste ; le présent document est
accompagné du \emph{Livre du Coffre}, complément opérationnel.

\textbf{D137} consigne la révision v0.6.1 : le niveau de sécurité d'un coffre neuf —
CF1, CF2 ou CF3 — devient l'arbitrage du développeur lors de la création de son projet
de développement. La règle D133 « un coffre neuf naît CF3 » est explicitement
invalidée et remplacée (§\ref{sec:evolutions061}) : imposer CF3 par défaut risque
d'être trop contraignant selon le projet. Le \code{init} de CF2 demeure un outil de
laboratoire (la spécification scellée CF2 fait foi) ; la naissance CF2 passe par
\code{init} CF1 puis la cérémonie \code{migrer} ; le coffre réel reste CF3 ; le
\emph{Livre du Coffre} passe à v1.1 dans le même geste.

Une décision ne se modifie jamais en silence : la révision remplace l'entrée par une note
d'invalidation explicite référençant la décision qui la supplante. La gouvernance
s'applique à elle-même — la présente révision en est la démonstration.

\appendix
\section{Empreintes md5 des modules du dépôt}\label{sec:empreintes}

Mesurées à la passation (commit de référence \code{c504436}, contenu canonique LF du
dépôt) — la génération de code se lit, ne se devine pas : un rapport de cérémonie
affiche les mêmes empreintes, le recoupement est mécanique.

\begin{table}[htbp]
\centering\footnotesize
\caption{Empreintes md5 des modules v06 (contenu canonique LF du dépôt, D133).}
\label{tab:empreintes}
\begin{tabularx}{\textwidth}{@{}l l Y@{}}
\toprule
\textbf{Module} & \textbf{md5} & \textbf{Rôle} \\
\midrule
\code{v06/coffre\_fort.py} (racine) & \code{8c79263b9b4694662fb722a67f14ee82} & coffre
CF1 (migration CF1 $\to$ CF2 ; terrain : divergence du clone à recouper) \\
\code{v06/coffre\_cf2.py} & \code{f50ce894db276818d68dd1e22a92fa4c} & coffre CF2 —
Shamir, sentinelle, porte D120 \\
\code{v06/coffre\_cf3.py} & \code{1cf18644d70073bf82f1f2847976f17d} & coffre CF3 —
serrure piégée, migration, destruction \\
\code{v06/sceau\_I.py} & \code{baf727058504490c288630dbc5bd9664} & sceau de
l'Intention — préparateur, dépositeur, vérificateur \\
\code{v06/ceremonie\_hors\_ligne.py} & \code{982bf0a402607b67a3960952a0adcb11} &
cérémonie air-gap (stdlib seule) \\
\code{v06/ceremonie\_gui.py} & \code{e4745580b19e7e5508f596b83bb1652f} & maître de
cérémonie graphique (v2.2, D131) \\
\code{v06/migration\_cf3\_gui.py} & \code{933e8405fa49ecad1d11861bd37d1c7a} & GUI de
migration CF2 $\to$ CF3 (v1.1, D132) \\
\code{v06/autotest\_cf3\_gui.py} & \code{bcd243448873d633852af27f204ecc27} & autotest
CF3 graphique \\
\code{v06/sentinelle\_cf2.py} & \code{bbc487423302c7f8f821ff348f3737c4} & sentinelle
CF2 (moteur laboratoire) \\
\code{v06/test\_cf2.py} & \code{f55fbee3391bbbb4466ba096c06b2ad7} & falsifiabilité
CF2 (Shamir, sentinelle, restauration) \\
\code{v06/test\_coffre\_cf3.py} & \code{33f460171f69aa59ae40192b5de46114} &
falsifiabilité CF3 (course de clôture, pannes) \\
\code{v06/test\_sceau\_I.py} & \code{a4b43fd5aa37acb37e98b77dca5747fc} &
falsifiabilité du sceau de l'Intention \\
\code{v06/test\_migration\_cf3.py} & \code{9e078dfdebea4bbe88b72fb71ab907bd} &
falsifiabilité de la migration CF2 $\to$ CF3 \\
\code{v06/test\_migration\_cf3\_gui.py} & \code{5f93800b887e44f52b1e998d5aabf160} &
falsifiabilité du GUI de migration \\
\code{v06/test\_ceremonie\_gui.py} & \code{e275efb83c104bcac30f05d96c7ce318} &
falsifiabilité du maître de cérémonie \\
\bottomrule
\end{tabularx}
\end{table}


Une empreinte divergente au rapport de terrain désigne soit une fin de ligne
différente (checkout Windows \code{autocrlf} : la commande git fait foi), soit une
génération de code périmée — le manifeste du kit tranche.

\section{Tableau décisionnel du coffre-fort (D75--D137)}\label{sec:annexedecisions}

Le journal complet vit dans \code{docs/DECISIONS.md} ; les entrées ci-dessous sont
celles qui touchent le coffre, sa garde et ses cérémonies — une ligne chacune, la
décision et ce qu'elle établit.

\begin{longtable}{@{}p{1.6cm} p{4.6cm} >{\raggedright\arraybackslash}p{9.5cm}@{}}
\toprule
\textbf{ID} & \textbf{Titre (abrégé)} & \textbf{Ce qu'elle établit} \\
\midrule
\endfirsthead
\toprule
\textbf{ID} & \textbf{Titre (abrégé)} & \textbf{Ce qu'elle établit} \\
\midrule
\endhead
\bottomrule
\endfoot
D75 & Coffre-fort des phrases maîtresses & conteneur TRIADE-CF1 : phrases par projet,
Argon2id/HKDF/AEAD, autotest de falsification, restauration post-réinitialisation. \\
D76 & Nettoyage du coffre Drive & douze livrables périmés archivés, miroir rendu fidèle
à git ; incident du reliquat consigné. \\
D80 & Rotation de lignée & la phrase exposée (D78) retirée du service, carte re-scellée
à faits inchangés, coffre roté. \\
D81 & Garde durable de la phrase de contrôle & \code{PHRASE-CONTROLE-COFFRE.txt} au
dépôt et au coffre Drive AVANT toute autre action — le fichier est la remise. \\
D84 & Licence BSL 1.1 & la barrière juridique de l'expropriation ; les preuves pour la
machine. \\
D85 & Attestation de non-exfiltration (axe 1) & quote TPM/SGX exigée du code manipulant
le coffre — exigence consignée, infrastructure sans TEE. \\
D86 & Interrupteur d'homme mort (axe 2) & clé maîtresse Shamir jamais entière ;
sentinelle ; destruction des fragments ; seuil T = 1 refusé. \\
D87 & Filigrane de l'Intention (axe 3) & arbres de couche I signés par la clé Ed25519
hors-ligne ; build et patch refusent l'intention non signée. \\
D105 & Axe 3 livré (laboratoire) & spec écrite-avant-code, cérémonie hors-ligne
autoportière, dix arbres, falsifiabilité 47/47 ; aucune clé réelle touchée. \\
D106 & Axe 2 livré (laboratoire) & coffre CF2 complet : Shamir (générateur primitif
0x03), sentinelle inotify + \code{/proc} double seuil, destruction trois passes + PV,
migration CF1 $\to$ CF2. \\
D107 & Maître de cérémonie graphique & un fichier, deux profils, zéro ligne de
commande. \\
D108 & Kit de cérémonie au Drive & \code{v06/} en mini-dépôt portable, MANIFESTE
d'empreintes, aucun secret. \\
D110 & La phrase de contrôle est à VOUS & alphanumérique, sans plafond — la discipline
D81 ne bouge pas. \\
D111--D118 & GUI de cérémonie, série & attachement nommable et gardé (D111), trois
chemins choisis dans la fenêtre (D112), boîte noire locale (D113), verdict ancré sans
répétition (D114), manifeste au contenu canonique LF (D115), fil de cérémonie par
preuves (D116), scénario avant tout (D117), page unique, dossier de cérémonie hors
dépôt, lanceurs (D118). \\
D119 & Dossier de cérémonie et fsync & le dossier naît avec la clé ; fsync de
répertoire POSIX-seul — le terrain Windows fait loi. \\
D120 & La sentinelle exigée au déverrouillage & la PORTE ferme le coffre quand le
témoin meurt ; ère sans moteur Windows : avertissement nommé à chaque déverrouillage. \\
D121 & Serrure à triple scellement (axe 4) & CF3 : le démon remplacé par un piège hors
machine ; cycle en deux actes ; erratum « 2 fragments » $\to$ TROIS. \\
D122 & Kit unique et trio systématique & une GUI, une boîte noire, un bouton « Envoyer
le rapport » pour chaque module à tester. \\
D123 & PV de destruction distincts & suffixe \code{-2} : deux destructions la même
seconde produisent DEUX fichiers (écrasement prouvé par le terrain D122). \\
D124 & Leçons GUI dans le berceau & \code{python -u}, tampons neutralisés,
\code{PYTHONIOENCODING} — les enfants ne mentent plus par buffer. \\
D125 & Preuve terrain CF3 (laboratoire) & la serrure CF3 prouvée avant tout outillage
de migration. \\
D126 & Re-cérémonie du sceau & accomplie par le développeur depuis son clone : dix
signatures vérifiées, commit poussé, la spec CF3 et l'erratum sous le sceau réel. \\
D127 & Outillage de migration CF2 $\to$ CF3 & \code{migrer} et \code{detruire-cf2},
GUI de migration, lanceurs ; re-lecture à froid avant remplacement atomique ;
destruction EXPLICITE après verdict humain. \\
D129 & Le pilier droit & fil d'étapes et journal dans un encart droit permanent, quel
que soit l'onglet. \\
D130 & Le kit ne dérive plus du miroir & \code{v06/} du kit déployé depuis
\code{git ls-files v06/} ; hook post-commit porte miroir PUIS kit. \\
D131 & Un bouton d'attachement par page & le doublon de l'accueil supprimé ; la note
fragments nomme les cinq originaux (\code{.keys/gardes-cf2/g1...g5}). \\
D132 & La grille du GUI de migration & chaque bouton « Parcourir » sur la ligne de son
champ ; confirmation de phrase étiquetée ; actes pleine largeur. \\
D133 & Passation v0.6.0 & le programme accompli passe en spécification ; la sentinelle
est \emph{par invocation} (la serrure) ; un coffre neuf naît CF3 ; M5 ouvert ; le
Livre du Coffre accompagne. \\
D137 & Arbitrage du niveau de naissance & le niveau d'un coffre neuf (CF1/CF2/CF3) est
choisi par le développeur à la création ; la règle « un coffre neuf naît CF3 » est
explicitement invalidée ; la chaîne de migration demeure disponible ; le coffre réel
reste CF3. \\
D160 & L'ère d'après M7 & l'ordre des familles est syscalls directs $\to$ TLS 1.3
pur $\to$ profils : la fondation I/O d'abord ; le jalon M8 naît sur le profil win
(domaine \code{sys}). \\
D162 & La cérémonie d'ère avant tout build & le garde-sceau NOMME la dette
(registre étendu non re-signé, arbre neuf) ; la table du sceau étendue à m8. \\
D163 & Continuité de clé par le déterminisme Ed25519 & re-signer un arbre inchangé
redonne exactement la signature déposée : la preuve de la réutilisation de la
clé, sans jamais voir la clé privée. \\
D165 & La finale d'une ère se scelle sur le recoup exigé comme pièce & le scellage
REFUSE sans le verdict du recoup vert ; le scellage est un fait de carte,
l'artefact ne bouge jamais. \\
D166 & La seconde famille ouverte : TLS 1.3 pur & le jalon M9, le domaine
\code{tls13}, l'ère 1.0.0 ; la dette de cérémonie nommée à la naissance. \\
D167 & Le chantier M9-(a) clos sur les 28 spécifications & la couche I complète
avant tout octet de notation ; les constantes du plan dérivées par arithmétique. \\
D168 & La charge de cérémonie se prépare au geste de la cérémonie & la charge de
l'ouverture était périmée par le geste des 24 specs ; l'ouverture annonce, la
cérémonie exige. \\
D169 & La consolidation v0.7.0 & le programme d'après M7 entre en spécification :
M8 franchi, M9 ouvert, le sceau à treize arbres, les cérémonies d'ère au
protocole ; la clause « hors périmètre courant » levée pour les deux premières
familles. \\
D170 & Le jalon M9 franchi : la finale 1.0.0 & le scellage sur le recoup terrain
exigé comme pièce (le miroir D165 à l'ère M9) ; un fait de carte, l'artefact ne
bouge jamais — \code{proxy.exe} byte-identique \code{446279f3…}. \\
D171 & Le test grandeur nature ouvert & « Loupe », l'éditeur hexadécimal au
panneau d'empreintes, arbitrée (contredisable au message suivant) ; le domaine
\code{loupe} fondé ; la dette de cérémonie nommée à la naissance ; le projet
fils COFFRE-FORT dégelé en parallèle. \\
D172 & Le chantier M10-(a) clos sur les 18 spécifications & la couche I complète
du TGN avant tout octet de notation (R2) ; le plan d'adressage arrêté avant la
notation (D161 rejoué). \\
D173 & Le chantier M10-(b) : le noyau noté-mesuré-prouvé & fichier + empreintes
(sha-256 par la carte, md5 la notation neuve RFC 1321) ; les héritages vivent à
leurs adresses scellées (R7) ; le mini-OS virtuel du harnais. \\
D174 & Le chantier M10-(c) : rendu + navigation & l'oracle à scénarios vert —
l'écran attendu octet par octet — et la falsification attrapée ; le mini-GUI
virtuel du harnais ; le plancher qui ne recule jamais. \\
D175 & Le chantier M10-(d), première moitié : l'édition & la famille édition
notée-mesurée-prouvée (le disque saisissable, la mort du harnais, le réseau de
la garde 16$\leftrightarrow$17$\leftrightarrow$18, l'écriture à l'octet prouvée
différentielle) ; le kit terrain attend la cérémonie d'ère. \\
D176 & La consolidation v0.8.0 & le programme d'après M9 entre en spécification :
M9 franchi (finale \code{1.0.0}), le TGN ouvert (« Loupe », la notation 18/18),
la dette du quatorzième arbre nommée — la cérémonie d'ère du TGN précède tout
build. \\
\end{longtable}